% !TeX root = ../测度论笔记.tex
% 由主文档通过 \input 载入；请编译主文档。
\setcounter{chapter}{4} % 保留原书章号。
\chapter{可测函数与简单函数逼近}
\label{chap:mt-functions}

\section{可测性是原像条件}
\begin{definition}
从可测空间 $(S,\mathcal A)$ 到 $(T,\mathcal D)$ 的映射 $f$ 称为\emph{可测映射（measurable map）}，若对每个 $B\in\mathcal D$，$f^{-1}(B)\in\mathcal A$。实值函数 $f:S\to\mathbb R$ 的可测性默认以 $\mathcal B(\mathbb R)$ 为目标集合族。
\end{definition}
可测性要求测量问题能够沿函数拉回底空间。$f^{-1}(B)=\{s:f(s)\in B\}$ 是原像，不要求 $f$ 有反函数，也不要求原像中只有一点。

\begin{proposition}[生成族判别]
实值函数 $f$ 可测，当且仅当对每个 $a\in\mathbb R$，集合 $\{f>a\}$ 可测。仅检查所有有理数 $a$ 也足够。
\label{prop:mt-measurable-test}
\end{proposition}
\begin{proof}
可测函数的射线原像必然可测。反向令 $\mathcal D_f=\{B\subseteq\mathbb R:f^{-1}(B)\in\mathcal A\}$，原像保持补集和可数并，所以 $\mathcal D_f$ 是 $\sigma$ 代数。若所有 $\{f>a\}$ 可测，则 $\mathcal D_f$ 包含生成 Borel $\sigma$ 代数的射线，从而包含所有 Borel 集。只知道有理阈值时，利用 $\{f>a\}=\bigcup_{q\in\mathbb Q,\ q>a}\{f>q\}$ 即可。
\end{proof}

连续实函数是 Borel 可测的，因为开集的原像开。$\mathbf1_A$ 可测当且仅当 $A$ 可测。两个可测映射的复合在目标、来源集合族相容时仍可测，因为 $(g\circ f)^{-1}(B)=f^{-1}(g^{-1}(B))$。

\section{运算与极限保持可测性}
若 $f,g$ 为有限实值可测函数，则 $f+g$、$fg$、$|f|$、$\max(f,g)$ 与 $\min(f,g)$ 均可测。例如
\[
 \{f+g>a\}=\bigcup_{q\in\mathbb Q}\bigl(\{f>q\}\cap\{g>a-q\}\bigr).
\]
平方可由连续函数复合处理，再利用 $fg=((f+g)^2-f^2-g^2)/2$ 处理乘法。扩展实值函数允许取 $\pm\infty$，其可测性仍可用实阈值判别，但不能对 $\infty-\infty$ 等未定义表达式直接做代数运算。

对可测函数列，$\sup_nf_n$ 可测，因为 $\{\sup_nf_n>a\}=\bigcup_n\{f_n>a\}$；取负号得到下确界的对应结论。因此
\[
 \liminf_nf_n=\sup_m\inf_{n\geq m}f_n,\qquad
 \limsup_nf_n=\inf_m\sup_{n\geq m}f_n
\]
均可测，逐点存在的极限也可测。若仅知道某函数与可测函数几乎处处相等，要保证它自身可测，还需要底空间完备，或另外明确要求该函数可测。

\section{简单函数的单调逼近}
\begin{definition}
取有限个实数值的可测函数称为\emph{简单函数（simple function）}，可以写成 $s=\sum_{j=1}^{m}a_j\mathbf1_{A_j}$，其中 $A_j$ 是两两不交的可测集。通过加入零值部分，可以使这些集合覆盖整个底空间。
\end{definition}

\begin{theorem}
对每个非负扩展实值可测函数 $f$，存在非负简单函数列 $s_n\uparrow f$。
\label{thm:mt-simple-approximation}
\end{theorem}
\begin{proof}
将值域按长度 $2^{-n}$ 分档，并在高度 $n$ 截断：
\[
 s_n(s)=
 \begin{cases}
  2^{-n}\lfloor 2^nf(s)\rfloor,&0\leq f(s)<n,\\
  n,&f(s)\geq n.
 \end{cases}
\]
其中 $\lfloor u\rfloor$ 表示不超过 $u$ 的最大整数。每个值对应的原像由可测水平集组成，所以 $s_n$ 是简单函数。分档逐次细化、截断高度上升，保证 $s_n\leq s_{n+1}$。对 $f(s)<\infty$，充分大的 $n$ 满足 $0\leq f(s)-s_n(s)<2^{-n}$；对 $f(s)=\infty$，有 $s_n(s)=n\to\infty$。
\end{proof}
这一定理是积分构造的关键。它不是把自变量轴切成等长区间，而是先按函数值分档，再测量每档所对应集合的大小。

\section{向量值映射、水平集与极限存在的集合}
当函数取值于 $\mathbb R^d$ 时，可测性可以逐坐标检查。设 $F=(f_1,\ldots,f_d)$，则 $F$ 关于 $\mathcal B(\mathbb R^d)$ 可测，当且仅当每个坐标函数可测。一个方向来自坐标投影连续；另一个方向则利用有理端点开矩形生成 $\mathbb R^d$ 的 Borel 集合族，而矩形的原像是有限个坐标原像的交。有限维的可数拓扑基使这一判别成为可能。

于是，若 $f_1,\ldots,f_d$ 可测且 $\Phi:\mathbb R^d\to\mathbb R$ 连续，则 $\Phi(f_1,\ldots,f_d)$ 可测。这统一解释了加法、乘法、绝对值、最大值以及许多连续变换为何保持可测性。若 $\Phi$ 仅为 Borel 可测，结论仍成立，因为可测映射的复合可测。应用时必须检查中间空间的集合族相容，不能只看两个映射分别被称为“可测”。

对函数列 $(f_n)$，收敛点组成的集合本身也可测。由实数的完备性，有限实值序列收敛当且仅当满足 Cauchy 条件。因此收敛点集合可以写成
\[
 C=\bigcap_{m=1}^{\infty}\bigcup_{N=1}^{\infty}
 \bigcap_{j,k\geq N}\{|f_j-f_k|\leq1/m\}.
\]
右侧只使用可数次集合运算，故可测。在 $C$ 上定义极限、在补集上定义为零，得到可测函数。这个构造在从依测度收敛提取子列，以及证明 $L^p$ 完备性时很有用：极限不必预先在全部点上存在，但可以在一个明确的可测集合上构造，再补成全空间函数。

\section{零测修改与可测版本}
若 $f$ 可测，$g=f$ 在某个可测零测集 $N$ 之外成立，而空间完备，则 $g$ 也可测。对每个阈值 $a$，集合 $\{g>a\}$ 与 $\{f>a\}$ 的对称差包含于 $N$；完备性保证这部分差异可测，因此阈值判别仍适用。若空间不完备，则可以取不可测的零测子集 $B\subseteq N$，令 $g=\mathbf1_B$，它与零函数几乎处处相等却不可测。这表明“几乎处处修改不影响可测性”含有一个不能省略的环境条件。

在实轴上，每个有限实值勒贝格可测函数都具有与它几乎处处相等的 Borel 可测版本。可先用简单函数逐点逼近；各简单函数的有限个水平集都与 Borel 集仅相差零测集，因此可替换成 Borel 可测简单函数，在必要时把重叠部分按固定次序剔除。全部替换产生的例外集合只有可数多个，其并仍为零测集。在这一并集之外，新旧简单函数列相同。新序列有限收敛的位置组成 Borel 集，在其上取极限、其余位置取零，就得到所需版本。

版本的概念以后还会出现在条件期望和随机过程里。它的含义是：当前理论只在几乎处处意义下确定对象，若要在每个点赋予具体数值，就需要选择一个代表。不能因为两个版本在零测集之外相同，就随意声称它们在某个指定点相等；也不能把不同参数下分别成立的几乎处处等式，未经检查就提升为对所有参数同时成立的等式。

\section{连续性与可测性的距离：Lusin 定理}
\label{sec:mt-lusin}
可测函数可以处处不连续，因而可测性远弱于连续性。但在有限测度区域内，舍去足够小的集合后，可测函数仍能表现出连续结构。下面的结果利用上一章的紧集逼近，把复杂水平集中的大部分质量保留在紧集内。它说明可测函数虽允许复杂局部行为，却并非完全无法由连续对象理解。

\begin{theorem}[Lusin 定理的有限测度版本]
设 $E\subseteq\mathbb R$ 勒贝格可测、$\lambda(E)<\infty$，$f:E\to\mathbb R$ 为有限实值可测函数。对每个 $\varepsilon>0$，存在紧集 $K\subseteq E$，使 $\lambda(E\setminus K)<\varepsilon$，并且限制函数 $f|_K$ 关于 $K$ 的子空间拓扑连续。
\label{thm:mt-lusin}
\end{theorem}
\begin{proof}
因为 $f$ 在每点有限，集合 $\{x\in E:|f(x)|>M\}$ 随 $M$ 增大递减到空集。有限测度下的从上连续性允许选择 $M$，使该集合测度小于 $\varepsilon/2$。记剩余集合为 $E_M$，在其上 $f$ 有界，故可由简单函数 $s_n$ 一致逼近。

对每个 $n$，把 $E_M$ 分成 $s_n$ 的有限个水平集，再利用紧集内逼近，从每个水平集中选紧子集，使所有遗漏测度之和小于 $\varepsilon2^{-n-1}$。这些紧子集的有限并记为 $K_n$。在 $K_n$ 上，$s_n$ 连续：不同的非空紧水平集两两不交，有限多个这样的紧集之间具有正距离，所以每个水平集在 $K_n$ 中既开又闭。

令 $K=\bigcap_nK_n$。它是紧集，且由次可加性，$\lambda(E\setminus K)<\varepsilon$。每个 $s_n|_K$ 连续，且仍一致收敛到 $f|_K$，故 $f|_K$ 连续。
\end{proof}

这里的结论是限制函数在 $K$ 上连续，不能理解为原函数在 $K$ 中每个点相对于整个实轴都连续。例如有理数示性函数处处不连续，但它在完全避开有理数的适当紧集上恒为零，限制函数当然连续。区分环境拓扑与子空间拓扑，就能解释这种表面上的矛盾。

Lusin 定理与 Egorov 定理关注不同问题。前者针对一个可测函数，舍去小集合后获得连续性；后者针对一个几乎处处收敛的函数列，舍去小集合后获得一致收敛。二者都把“几乎处处”的信息转化为大集合上的更强性质，但所用结论和所控制的对象不同，不能互相替代。

\section{习题}
\begin{enumerate}
 \item 证明 $\{f=g\}$ 是可测集，其中 $f,g$ 是有限实值可测函数。
 \item 证明非负简单函数的上述逼近对有界 $f$ 最终满足一致误差小于 $2^{-n}$。
 \item 解释在一个不完备测度空间中，零函数的零测集修改为什么可能不可测。
\end{enumerate}

